\chapter{แหล่งเรียนรู้และบรรณานุกรมเพิ่มเติม}
\label{app:D}

\section*{ง.1 เอกสารมาตรฐานและข้อกำหนด}

\begin{itemize}
    \item \textbf{Khronos Group} (2023). \textit{OpenXR 1.1 Specification}. The Khronos Group Inc. \\
          เข้าถึงได้จาก: \url{https://www.khronos.org/registry/OpenXR/specs/1.1/}
    \item \textbf{Khronos Group} (2023). \textit{OpenXR Reference Guide}. The Khronos Group Inc. \\
          เข้าถึงได้จาก: \url{https://www.khronos.org/files/openxr-10-reference-guide.pdf}
    \item \textbf{W3C Immersive Web Working Group} (2024). \textit{WebXR Device API}. \\
          เข้าถึงได้จาก: \url{https://www.w3.org/TR/webxr/}
\end{itemize}

\section*{ง.2 ซอฟต์แวร์และเครื่องมือโอเพนซอร์ส}

\begin{table}[H]
\centering
\caption{เครื่องมือและไลบรารีที่ใช้ในการพัฒนา OpenXR STEM}
\begin{tabularx}{\textwidth}{lYY}
\toprule
\textbf{เครื่องมือ} & \textbf{คำอธิบาย} & \textbf{URL} \\
\midrule
Monado & OpenXR Runtime (Linux) & gitlab.freedesktop.org/monado \\
PyOpenXR & Python binding สำหรับ OpenXR & pypi.org/project/pyopenxr \\
MediaPipe & Hand tracking และ AI pipeline & mediapipe.dev \\
NumPy & Scientific computing & numpy.org \\
Three.js & WebXR 3D Graphics & threejs.org \\
A-Frame & WebXR Framework & aframe.io \\
\bottomrule
\end{tabularx}
\end{table}

\section*{ง.3 แหล่งเรียนรู้ออนไลน์ที่แนะนำ}

\begin{itemize}
    \item \textbf{Khronos OpenXR Tutorial} — บทเรียน C++ OpenXR ระดับเริ่มต้นถึงขั้นสูง \\
          \url{https://openxr-tutorial.com}
    \item \textbf{Monado Developer Documentation} — คู่มือการพัฒนาและ API reference \\
          \url{https://monado.freedesktop.org/}
    \item \textbf{WebXR Samples} — ตัวอย่างการพัฒนา WebXR บน Web Browser \\
          \url{https://immersive-web.github.io/webxr-samples/}
    \item \textbf{MediaPipe Solutions} — คู่มือการใช้งาน Hand Landmark Detection \\
          \url{https://ai.google.dev/edge/mediapipe/solutions/vision/hand\_landmarker}
\end{itemize}

\section*{ง.4 วารสารวิชาการที่เกี่ยวข้อง}

\begin{itemize}
    \item \textit{IEEE Transactions on Visualization and Computer Graphics} (TVCG)
    \item \textit{Computers \& Education} — Elsevier
    \item \textit{International Journal of Human-Computer Studies} — Elsevier
    \item \textit{Journal of Educational Technology \& Society} — JSTOR
    \item \textit{วารสารวิทยาศาสตร์และเทคโนโลยี} มหาวิทยาลัยราชภัฏ
\end{itemize}
